\section{Entering the Feast}\label{sec:entering}

Christmas begins before midnight has struck. At First Vespers on the evening of
24~December, the Church lays aside the violet vestments of Advent, clothes her
altars in white, and turns from expectation to adoration. While the world around
her often empties its streets on 26~December and strips its trees before the new
year, the 1962 Roman liturgy refuses to hurry past Bethlehem. Twelve full days
of rejoicing---from First Vespers of the Nativity on 24~December through None of
5~January---belong to the feast of the Word made flesh.

If you are opening Christmastide on the afternoon of 24~December, take a quiet
quarter-hour before evening prayer to settle a simple rule of daily Scripture
and prayer that can survive holiday travel, guests, and winter errands
(Section~\ref{sec:prayer-home}). Spend a few minutes with the mystery of the
wonderful exchange and the ancient Roman stations of the three Christmas Masses,
so that when the priest intones the \textit{Gloria in excelsis} at midnight, you
stand inside the prayer the Church has kept for sixteen centuries.

If you are arriving late---on the evening of 28~December, on New Year's morning,
or during the quiet weekdays of early January, tempted to suppose that the feast
has already slipped past you---do not wait for next winter. Turn directly to
today's calendar date in Appendix~\ref{app:daily-course}, read the appointed
Scripture and the Fathers for that day, and pray the day's Collect before sleep.
As St.~Bernard of Clairvaux told his monks in the Octave, the Church sets a
whole train of solemnities directly after the Nativity precisely to keep
devotion alive and show forth the fruit of Christ's coming, ``by representing
these three festivals as a continuation, so to speak, of Christmas.'' Even a
single day kept beside the crib belongs to the whole feast.

In the 1960 \textit{Codex Rubricarum} that governs the 1962 Roman Rite, the
larger Christmas cycle (\textit{Tempus Natalicium}) runs from First Vespers of
24~December through 13~January inclusive, and divides cleanly into two distinct
seasons. The first is \textbf{Christmastide} (\textit{Tempus Nativitatis},
General Rubric~72~a), extending from First Vespers of 24~December through None
of 5~January inclusive. The second is \textbf{Epiphanytide} (\textit{Tempus
Epiphaniae}, General Rubric~72~b), opening at First Vespers of the Epiphany on
the evening of 5~January and closing on 13~January. Christmastide holds that
crisp twelve-day arc from the evening of 24~December right to the threshold of
the Epiphany star on 5~January.

Worshipping in the 1962 Roman Rite draws you into an ordered rhythm of Mass and
Hours whose landmarks quickly grow familiar. At the center stands the
\textbf{Mass}. On 25~December the Missal appoints not one Mass but three---at
Midnight, at Dawn, and in the Day---each with its own ancient Roman
\textbf{station} (\textit{Statio}), chants, readings, and prayers. Listen closely
to the \textbf{Collect} (\textit{Oratio}) before the Epistle, where the priest
gathers the whole mystery of the day into one sentence, and to the
\textbf{Preface of the Nativity} and the proper \textbf{\textit{Communicantes}}
inside the Roman Canon, which bind the sacrifice of the altar directly to the
night and day of Bethlehem through 1~January.

Immediately around the crib stand the \textbf{Octave of the Nativity}
(25~December through 1~January)---one of only three octaves retained by the
1962 books alongside Easter and Pentecost---and the saints who surround the
newborn King: St.~Stephen (26~December), St.~John the Evangelist (27~December),
and the Holy Innocents (28~December), followed by the remaining Octave days
carrying commemorations of St.~Thomas of Canterbury (29~December) and
St.~Sylvester~I (31~December), the \textbf{Sunday within the Octave}, the
first-class \textbf{Octave Day of the Nativity} on 1~January, the \textbf{Feast
of the Most Holy Name of Jesus} (kept on the Sunday between 2 and 5~January, or
on 2~January when no Sunday falls there), and the quiet weekdays and Marian
Saturdays of 2--5~January, closing with the commemoration of Pope
St.~Telesphorus on 5~January.

Surrounding the Mass turns the \textbf{Divine Office}: \textbf{Matins} in the
night watches, with Isaiah, the letters of St.~Paul, and sermons of the Fathers;
\textbf{Lauds} at dawn, singing \textit{A solis ortus cardine}, ``From where the
sunrise had its birth''; the daytime Hours of \textbf{Prime}, \textbf{Terce},
\textbf{Sext}, and \textbf{None}; \textbf{Vespers} at dusk, with the hymn
\textit{Jesu Redemptor omnium}, ``Jesus, the Ransomer of man,'' and Mary's
\textit{Magnificat}; and \textbf{Compline} before sleep, where the Marian
antiphon \textit{Alma Redemptoris Mater} turns at First Vespers of Christmas to
its Christmastide versicle---\textit{Post partum, Virgo, inviolata permansisti},
``After thy delivery thou still remainest a Virgin undefiled''---and the Collect
\textit{Deus, qui salutis aeternae}. Lay faithful are not bound to the clerical
Office, yet praying even one Hour or saying the Collects, antiphons, and hymns
in Appendix~\ref{app:prayers} joins a household to the praise of the whole
Church.

Why does the Church return to these same readings and melodies every winter?
Pope Pius~XI gave the answer in his encyclical \textit{Quas primas} (1925): for
imbuing the faithful with the truths of faith and lifting them to inward joy,
the annual celebrations of the sacred mysteries carry far greater force than
even the weightiest written decrees. Documents speak once and reach mostly the
learned; feasts speak every year in perpetuity to every believer, forming both
mind and heart---the whole human person. G.~K.~Chesterton waved away modern
impatience with recurring holy days in a single tart question in
\textit{Orthodoxy} (1908): ``How can anything be up to date? a date has no
character. How can one say that Christmas celebrations are not suitable to the
twenty-fifth of a month?''

On 25~December~380, in Constantinople, St.~Gregory of Nazianzus opened his
thirty-eighth Oration with the summons that still greets every Christian
entering the feast:

\begin{quote}
Christ is born, glorify ye Him. Christ from heaven, go ye out to meet Him.
Christ on earth; be ye exalted. Sing unto the Lord all the whole earth; and
that I may join both in one word, Let the heavens rejoice, and let the earth be
glad, for Him Who is of heaven and then of earth. Christ in the flesh, rejoice
with trembling and with joy; with trembling because of your sins, with joy
because of your hope. Christ of a Virgin; O ye Matrons live as Virgins, that ye
may be Mothers of Christ.
\end{quote}

Throughout the pages ahead, psalms carry their Latin Vulgate numbering first,
with Hebrew numbering in brackets where the two differ---for example,
Psalm~109~[110] or Psalm~97~[98]. Sources, translations, and calendar rules are
recorded in Appendix~\ref{app:about} and Appendix~\ref{app:references}.
